% abstract.tex -- single-study version, 2026-10-31 (Shlok approved: one study, abstract without numbers; PROFESSOR-SUGGESTIONS-ANALYSIS-2026-10-31.md).
% No number, interval, p-value, sample size or parameter count; every number stays in Results. One paragraph, no abbreviation, symbol,
% equation or citation. Deviation from the A6.12 sentence templates (numbers moved to Results) is listed in the Deviations section.
% Claims and their sources (all in Results, with run comments there):
% run: runs/P2-LRF-L1/a6_analysis.json:family_L.L1.label (= confirmed); family_L.L-LEN.label (= rejected); family_L.L-NI.{non_inferior (= true), margin_note}
% run: runs/P2-LRF-DN/dn_analysis.json:cells.{noise0.1,jpeg10}.arms.{DN+NONE,DN+LRF}.vs_LRF (denoise before scoring: above LRF under noise; DN+LRF below LRF under JPEG; DN+NONE under JPEG interval contains 0)
% run: runs/P2-LRF-ABL-SCORE/ablation_analysis.json:rows[id=A-LIN] (linear logit: noise interval contains 0; at or above LRF in the other three conditions)
Image-text verification, the choice of the true caption of an image from two candidates, becomes harder when the image is damaged.
We study this problem inside one frozen pretrained unified discrete diffusion model, UniDisc, which gives every observed image token a
probability given a caption, with no external encoder.
For each token, the log-probability ratio between the true and a wrong caption is evidence, and the standard rule adds the evidence of
all tokens with equal weight.
We learn a per-token reliability weighting of this evidence from features computed without the caption, fit it on design images, and
freeze it.
A preregistered one-sided exact paired test on fresh Common Objects in Context validation images under Gaussian noise at one level
confirmed that the frozen weighting raised verification accuracy, including when the wrong caption was matched to the true one in length.
On clean images the weighting was non-inferior within a preregistered margin; the margin was fixed after an exploratory clean-image
interval had been seen.
Two exploratory comparisons qualify this result.
A classical pixel denoiser applied before scoring reached higher accuracy than the weighting under Gaussian noise, and combined with the
weighting it reached lower accuracy under Joint Photographic Experts Group compression.
A linear logit on the same features reached similar accuracy under Gaussian noise and higher accuracy under compression.
An exploratory sweep on the same images found the gain at milder and stronger Gaussian noise as well.
The confirmatory evidence covers one model, one dataset, and one noise level; it does not show that the learned weights track token
reliability, or that the gain holds for other models or corruptions.
